\documentclass[10pt]{article}
\usepackage[T1]{fontenc}
\usepackage[utf8]{inputenc}
\usepackage[french]{babel}
\usepackage{amsmath,amssymb,mathrsfs,mathpazo}
\usepackage{geometry}
\geometry{a4paper,margin=2cm}
\usepackage{graphicx,array,tabularx,tabularray,colortbl,multirow,multicol,booktabs,enumitem,pifont,ulem,fancybox,diagbox,cancel}
\usepackage{pstricks,pstricks-add,pst-plot,pst-tree,pst-node,pst-text,pst-func,pst-math,pst-bspline,pst-3dplot,pst-eucl,pst-solides3d}
\let\pscalculate\undefined
\usepackage{pst-bezier}
\usepackage{tikz,pgfplots,tkz-tab}
\usetikzlibrary{arrows,arrows.meta,patterns,calc,positioning,trees,decorations.markings,shapes.geometric,shapes.misc,shapes.arrows,intersections,decorations.pathmorphing,decorations.pathreplacing,backgrounds,shadows,babel}
\pgfplotsset{compat=1.18}
\usepgfplotslibrary{fillbetween}
\usepackage[np]{numprint}
\usepackage{listings,makecell,pas-tableur,twoopt}
\let\barre\undefined
\usepackage[pstricks]{tablvar}
\definecolor{darkgreen}{RGB}{0,100,0}
\providecommand{\red}{\color{red}}
\providecommand{\blue}{\color{blue}}
\providecommand{\green}{\color{green}}
\providecommand{\black}{\color{black}}
\providecommand{\white}{\color{white}}
\providecommand{\gray}{\color{gray}}
\providecommand{\cyan}{\color{cyan}}
\providecommand{\magenta}{\color{magenta}}
\providecommand{\yellow}{\color{yellow}}
\providecommand{\violet}{\color{violet}}
\providecommand{\orange}{\color{orange}}
\providecommand{\euro}{\texteuro}
\providecommand{\EUR}{\texteuro}
\providecommand{\ieme}{\textsuperscript{e}}
\providecommand{\degres}{\ensuremath{^\circ}}
\providecommand{\degree}{\ensuremath{^\circ}}
\providecommand{\vv}[1]{\overrightarrow{#1}}
\providecommand{\Coord}[1]{#1}
\providecommand{\decosix}{}
\providecommand{\starredbullet}{\ensuremath{\bullet}}
\pagestyle{empty}
\setlength{\parindent}{0pt}
\usepackage{adjustbox}\DeclareUnicodeCharacter{2236}{:}\DeclareUnicodeCharacter{2190}{\ensuremath{\leftarrow}}
\providecommand{\R}{\mathbb{R}}
\providecommand{\N}{\mathbb{N}}
\providecommand{\Z}{\mathbb{Z}}
\providecommand{\Q}{\mathbb{Q}}
\providecommand{\D}{\mathbb{D}}
\providecommand{\C}{\mathbb{C}}
\providecommand{\e}{{\mathrm{e}}}
\providecommand{\vect}[1]{\overrightarrow{#1}}
\providecommand{\vectt}[1]{\overrightarrow{\mathrm{#1}}}
\providecommand{\vvt}[1]{\overrightarrow{\mathrm{#1}}}
\providecommand{\np}[2][]{{#2~#1}}
\providecommand{\vv}[1]{{\overrightarrow{#1}}}
\renewcommand{\d}{{\mathrm{d}}}
\providecommand{\degres}{{^\circ}}
\providecommand{\degre}{{^\circ}}
\providecommand{\blue}{\color{blue}}
\providecommand{\red}{\color{red}}
\providecommand{\enskip}{\quad}
\providecommand{\llbracket}{\lbrack\!\lbrack}
\providecommand{\rrbracket}{\rbrack\!\rbrack}
\providecommand{\euro}{{\text{€}}}
\providecommand{\no}{{\text{n°}}}
\providecommand{\interval}[2]{{[#1;#2]}}
\providecommand{\intervFO}[2]{{[#1;#2[}}
\providecommand{\dotfill}{\ldots}

\providecommand{\ds}{\displaystyle}
\providecommand{\textding}[1]{#1}
\providecommand{\Oij}{\(\mathrm{O};\vec{\imath},\vec{\jmath}\)}
\providecommand{\Oijk}{\(\mathrm{O};\vec{\imath},\vec{\jmath},\vec{k}\)}
\providecommand{\Ouv}{\(\mathrm{O};\vec{u},\vec{v}\)}
\providecommand{\intervFF}[2]{\ensuremath{[#1;#2]}}
\renewcommand{\euro}{\ifmmode\text{\texteuro}\else\texteuro\fi}
\DeclareUnicodeCharacter{00B0}{\ifmmode^{\circ}\else\textdegree\fi}
\DeclareUnicodeCharacter{20AC}{\ifmmode\text{\texteuro}\else\texteuro\fi}
% Énoncés : Éducation nationale. Transcription et corrigés : APMEP, Denis Vergès et contributeurs. Découpage et classement : CoopMaths, Sébastien Lozano, Éric Elter et contributeurs. Adaptation de la présentation : Maths Explorer.
% CC BY-SA : https://coopmaths.fr/www/about/
% https://coopmaths.fr/alea/static/bac/2024/tex/bac_2024_06_sujet2_polynesie_1_cor.tex
\begin{document}
\noindent\textit{APMEP / CoopMaths. CC BY-SA. Présentation adaptée par Maths Explorer.}\par\medskip
\medskip





















\medskip

\begin{enumerate}
	\item 
	

D'après l'énoncé $P(J) = 0,6$ et $P_J(S) = \dfrac89$.

On a donc $P(J \cap S) = P(J) \times P_J(S) = 0,6 \times \dfrac89 = \dfrac{4,8}{9} = \dfrac{3 \times 1,6}{3 \times 3} = \dfrac{1,6}{3}  = \dfrac{5 \times 1,6}{5 \times 3}  = \dfrac{8}{15}$.
\end{enumerate}



\begin{enumerate}[resume]
\item 
	\begin{enumerate}
		\item 
On sait que $P(S) = \dfrac23$.

On commence l'arbre de probabilités pondéré suivant :

\begin{center}
\pstree[treemode=R,levelsep=3cm]{\TR{}}
{\pstree{\TR{$J~$}\taput{0,6}}
	{
	\TR{$S$} \taput{$\frac89$}
	\TR{$\overline{S}$} \tbput{$\frac19$}
	}
\pstree{\TR{$\overline{J~}$}\tbput{0,4}}
	{\TR{$S$} \taput{$x$}
	\TR{$\overline{S}$} \tbput{$1 - x$}
	}
}
\end{center}

D'après la loi des probabilités totales :

$P(S) = P(J \cap S) + P\left(\overline{J} \cap S \right)$, soit $\dfrac23 = \dfrac{8}{15} + 
P\left(\overline{J} \cap S \right) \iff P\left(\overline{J} \cap S\right)  = \dfrac23 -  \dfrac{8}{15} =$

$\dfrac{10}{15} - \dfrac{8}{15} = \dfrac{2}{15}$.

La probabilité que la personne choisie n'a pas l'intention de regarder les JO à la télévision et a une activité sportive régulière est égale à $\dfrac{2}{15}$;
		\item 
		
On a  $P_{\overline{J}}(S) = \dfrac{P\left(\overline{J} \cap S\right)}{P\left(\overline{J}\right)} = \dfrac{\frac{2}{15}}{0,4} = \dfrac{\frac{2}{15}}{\frac{2}{5}}  = \dfrac{2}{15} \times \dfrac{5}{2} = \dfrac13 = x$.
	\end{enumerate}
\end{enumerate}

\medskip



\begin{enumerate}[resume]
\item 
	

	

	\begin{enumerate}
		\item 
Les personnes sont choisies au hasard et chacune d'elles a une pratique sportive régulière avec une probabilité de $\dfrac23$.

La variable aléatoire $X$ suit donc une loi binomiale de paramètres $n = 30$ et 

$p = \dfrac23$.
		\item 
		On a $P(X = 16) = \displaystyle\binom{30}{16} \left(\dfrac23\right)^{16}\times \left(1 - \dfrac23\right)^{30 - 16} = \displaystyle\binom{30}{16} \left(\dfrac23\right)^{16}\times \left(\dfrac13\right)^{14} \approx 0,0462$ soit 0,046 au millième près à la calculatrice.
		\item 
		

		

On a $\dfrac{\np{10000}}{380}\approx 26,3$ : on peut donc offrir 26 entrées gratuites au maximum.

La calculatrice donne $P(X \leqslant 26) \approx 0,9967$ soit environ 0,997.
La probabilité que le budget soit insuffisant est donc égale à environ $1 - 0,997$ soit environ 0,003 ou trois millièmes.
	\end{enumerate}
\end{enumerate}

\bigskip

\end{document}
